% V12 reader-first revision. V11 experimental and mathematical records are preserved.
\newcommand{\HalfMoonIntroduction}{%
\section{Introduction}\label{sec:introduction}
Direction normalization is useful when orientation matters more than magnitude: it suppresses variation in distance from a reference origin. That same operation can destroy class information when distance matters. The half-moon task makes the conflict concrete. Its two classes occupy thick semicircular bands---an upper half-ring and a translated lower half-ring---whose positions determine whether they are disjoint or overlap (Fig.~\ref{fig:geometry}). We establish that direction alone loses class information at every origin for specified arrangements and that one near/far radius bit restores the decisions in a disjoint case.

At vertical displacement $h=0.5$, take the midpoint origin $c=(0.5,0.25)$. The class-0 point $x_0=(1,0.25)$ and class-1 point $x_1=(2,0.25)$ both lie inside their respective half-rings. Each lies directly to the right of $c$, so their normalized directions coincide at $(1,0)$ despite opposite labels. Their distances from $c$ are $0.5$ and $1.5$; a near/far threshold of $0.8$ distinguishes them. This constructed example explains one lost decision, not a positive population error. The proof below establishes conflicting decisions on a set of directions with positive measure, and does so for every origin.

We quantify the loss with \emph{Bayes risk}: the lowest possible misclassification probability when the class distributions are known and the classifier makes the best decision permitted by its input. Comparing this ideal error for raw coordinates with that for directions isolates information removed by preprocessing. The difference is distinct from the extra error of a trained network using finite data and an imperfectly learned decision rule. It also defines a harder question than whether a few chosen origins give poor accuracy: can \emph{any} origin, including one arbitrarily far away, remove the representation-induced error?

Spatial-sign preprocessing and geometric classifiers already use normalized directions \citep{Serneels2006,Graf2003}. Decision-sufficiency theory establishes when a transformation retains optimal class decisions \citep{Silva2012,Silva2024,Gyorfi2023}, while recent \emph{Pattern Recognition Letters} studies connect classification performance to feature geometry and normalization \citep{Poblete2026,Ahmadian2025}. The unresolved geometric question for the present populations is whether these decisions can be preserved by moving a shared origin, and what restricted radial information would recover them when they cannot.

We answer it by deriving the probability of opposing decisions merged along each ray, proving a uniform distant-origin limit to orthogonal-projection risk, and constructing explicit positive lower bounds over the whole plane. A common-threshold radius bit then selects between two angular decisions and exactly recovers the labels in one disjoint half-moon configuration; the admissible threshold interval also determines its radius-error tolerance. Controlled neural experiments test how close learned decisions come to the corresponding population risks, using 500 candidate fits and 250 selected evaluations across the main and repair studies.
}

\newcommand{\HalfMoonBackground}{%
\section{Background and related work}\label{sec:background}
\paragraph{Representations and achievable classification error.}
A preprocessing map determines what information a classifier can use. The spatial-sign map of an observation $x$ about a shared origin $c$ is $T_c(x)=(x-c)/\|x-c\|$ for $x\ne c$: it reports a unit direction and ignores the radius $\rho=\|x-c\|$ \citep{Serneels2006}. Direction and radius together recover the original point by $x=c+\rho T_c(x)$. Retaining both therefore supplies a full-information control. A Bayes classifier selects the most probable label given the supplied input. Its Bayes risk is the residual probability of error even under that optimal rule. A larger Bayes risk after normalization means that the representation itself prevents some correct decisions; a neural model cannot recover them by more training.

\paragraph{What earlier geometric and normalization studies establish.}
Spatial-sign preprocessing and normalized-feature classification establish ways to compare directions \citep{Serneels2006,Graf2003}. In \emph{Pattern Recognition Letters}, Poblete et al.~\citep{Poblete2026} compare classifier rules on fixed learned features across eight datasets: rankings vary with geometry, while the predicted directional association of radial compatibility is unsupported. Their evidence concerns choosing a rule for an existing representation. Ahmadian and Lindsten~\citep{Ahmadian2025} evaluate activity novelty detection with radial kernels, normalization and quantization; their ablations report different performance for normalized and unnormalized variants. These observations motivate a sharper question: before fitting any classifier, can a change of normalization origin prevent the loss of optimal class decisions?

\paragraph{Decision preservation and radial information.}
Silva and Narayanan~\citep{Silva2012} analyze representation quality through Bayes and learned error; later decision-sufficiency results distinguish preserving a label distribution from preserving the decision that minimizes a given loss \citep{Silva2024,Gyorfi2023,Silva2025}. Dubois et al.~\citep{Dubois2021} identify unit direction as an invariant of positive scaling. The key condition is intuitive: points collapsed to the same direction must require the same optimal decision for direction alone to be sufficient. Our half-moons supply a continuous geometry in which that condition can fail and in which the normalization origin itself varies.

A representation can instead preserve only the extra information needed for the decision. Bayes-aware quantization and learned cut points investigate compressed decision inputs \citep{OehlerGray1995,Senozan2024}. Our repair asks a precise question: can \emph{one shared near/far radius threshold}, combined with the retained direction, preserve decisions across both continuous class regions? The threshold then acts as a minimal decision cue rather than a measurement that reconstructs the full observation.

\paragraph{Why all origins require a separate analysis.}
Projection methods study the error incurred by a selected lower-dimensional view \citep{Hamsici2007,Hamsici2008,Oh2025}. From a very distant origin, rays crossing a bounded population become nearly parallel, and normalized direction approaches an orthogonal projection in its decision-making effect \citep{Orponen2019}. Thus proving positive loss for each finite origin alone does not settle whether the infimum over \emph{all} origins is positive. Finite planar separability geometry provides another ingredient: connected regions whose convex hulls intrude into each other cannot always be isolated into different directional sectors \citep{Hurtado2001}. We combine these geometric observations with population risk to establish both finite-origin loss and its strictly positive global floor, then ask how a restricted radius code changes that conclusion.
}

\newcommand{\HalfMoonDiscussion}{%
\section{Discussion}\label{sec:discussion}
The experiments distinguish error already present in the original coordinates, additional error forced by direction-only preprocessing, and excess error from fitting a network. In the disjoint midpoint condition, the small learned excess over the Bayes floor (Table~\ref{tab:neural-main}) shows that most direction-only errors arise from the input representation. In the overlapping condition, raw-coordinate ambiguity already exists and normalization adds further loss. Accuracy comparisons alone would not reveal these two mechanisms.

Poblete et al.~\citep{Poblete2026} show that classifier rankings depend on the geometry of fixed learned features; our all-origin criterion instead determines when \emph{no} classifier can recover decisions removed by direction normalization. Ahmadian and Lindsten~\citep{Ahmadian2025} document normalization-sensitive novelty-detection performance; our Bayes-risk analysis isolates representation loss before learning. Relative to decision-preservation theory \citep{Silva2012,Silva2024,Gyorfi2023}, the ray-cancellation identity and distant-origin limit give a geometric test for a whole family of normalization maps, including origins receding to infinity.

The one-bit result specifies the information needed to repair the disjoint midpoint geometry. Radius alone has $50\%$ Bayes error under the equal-prior reflected sampling laws, yet its near/far bit selects between two angular separators when combined with direction. The certified threshold interval applies to the explicit decoder; the independent neural evaluation shows how closely a trained model uses that representation, without equating its few residual errors with the decoder's exact guarantee. The resulting evaluation sequence is to identify the retained class information, calculate its optimal risk, and then assess learning error relative to that risk.
}
